% =====================================================================
% Problem Formulation
% Cross-Version Adapter Portability via Representational Alignment
% =====================================================================
\section{Problem Formulation}
\label{sec:problem}

\subsection{Setup and Notation}

Let $\mathcal{F}$ denote a model \emph{family} (e.g.\ Gemma, Qwen) and let
$M_v$ and $M_{v+1}$ be two \emph{base} (pre-trained, instruction-tuned, or
foundation) checkpoints from $\mathcal{F}$ released at successive versions
$v$ and $v{+}1$. The two checkpoints typically share architecture and
tokenizer but differ in weights due to retraining; we make no assumption
that they share a representation space.

A \emph{LoRA adapter}~\citep{hu2022lora} $A$ trained on top of $M_v$ is a set
of low-rank updates over a collection of adapted linear modules
$\ell \in \mathcal{L}$. Module $\ell$ realizes a linear map with frozen
weight $W_\ell \in \mathbb{R}^{d^{\mathrm{out}}_\ell \times d^{\mathrm{in}}_\ell}$,
and LoRA augments it with
\begin{equation}
  \Delta W_\ell \;=\; \tfrac{\alpha}{r}\, B_\ell A_\ell,
  \qquad
  B_\ell \in \mathbb{R}^{d^{\mathrm{out}}_\ell \times r},\;
  A_\ell \in \mathbb{R}^{r \times d^{\mathrm{in}}_\ell},
\label{eq:lora}
\end{equation}
with rank $r \ll \min(d^{\mathrm{in}}_\ell, d^{\mathrm{out}}_\ell)$ and scale
$\alpha$. The fine-tuned model is denoted $M_v \oplus A$, i.e.\ each adapted
module computes $y = (W_\ell + \Delta W_\ell)\,x$.

The adapter $A$ is trained on a task dataset $\mathcal{D}_{\mathrm{task}}$
(potentially private, expensive, or no longer available). We additionally
assume access to a small, \emph{unlabeled} \emph{anchor} set
$\mathcal{D}_{\mathrm{anchor}}$ ($|\mathcal{D}_{\mathrm{anchor}}| \ll
|\mathcal{D}_{\mathrm{task}}|$, e.g.\ a few thousand generic text sequences)
that can be passed through both base checkpoints.

\subsection{Objective}

Let $\mathcal{P}(M \oplus A)$ be task performance of model $M$ equipped with
adapter $A$. We seek a transformation $T$ that transplants $A$ from $M_v$ onto
$M_{v+1}$ \emph{without access to $\mathcal{D}_{\mathrm{task}}$}:
\begin{equation}
  T^\star \;=\;
  \arg\max_{T \in \mathcal{T}}\;
  \mathcal{P}\!\left(M_{v+1} \oplus T(A)\right),
  \qquad
  \text{using only } M_v,\, M_{v+1},\, A,\, \mathcal{D}_{\mathrm{anchor}}.
\label{eq:objective}
\end{equation}
We compare against two reference points: the \emph{naive transfer} baseline
$T = \mathrm{id}$ (drop $A$ onto $M_{v+1}$ unchanged), and the \emph{oracle}
$A^\star = \arg\max_A \mathcal{P}(M_{v+1}\oplus A)$ obtained by retraining on
$\mathcal{D}_{\mathrm{task}}$. A successful $T$ closes the gap between naive
transfer and the oracle at a small fraction of the oracle's compute, and
without task labels.

\subsection{Why Naive Transfer Fails: Representational Drift}

Equation~\eqref{eq:lora} shows the adapter acts on the module's \emph{input}
activation and writes into its \emph{output} activation. Concretely, for input
$x$ in the input representation space of $M_v$, the adapter contributes
$\tfrac{\alpha}{r} B_\ell (A_\ell x)$. The adapter is therefore \emph{bound}
to the geometry of $M_v$'s activation spaces: $A_\ell$ reads coordinates in
$M_v$'s input basis and $B_\ell$ writes coordinates in $M_v$'s output basis.

When the same module's activation geometry differs in $M_{v+1}$, the directions
$A_\ell$ reads from and $B_\ell$ writes to no longer correspond to the same
semantic features, so the update is applied in the wrong basis. We quantify
this \emph{representational drift} per module with linear Centered Kernel
Alignment~\citep{kornblith2019cka}: collecting input activation matrices
$H^{v}_\ell, H^{v+1}_\ell \in \mathbb{R}^{N \times d^{\mathrm{in}}_\ell}$ over
$\mathcal{D}_{\mathrm{anchor}}$ ($N$ tokens),
\begin{equation}
  \mathrm{drift}_\ell \;=\; 1 - \mathrm{CKA}\!\left(H^{v}_\ell, H^{v+1}_\ell\right),
  \qquad
  \mathrm{CKA}(X,Y) = \frac{\lVert Y^\top X \rVert_F^2}
       {\lVert X^\top X \rVert_F \,\lVert Y^\top Y \rVert_F}.
\label{eq:cka}
\end{equation}
A central empirical claim of this paper (\S\ref{sec:drift}) is that
$\mathrm{drift}_\ell > 0$ across versions and grows with depth, and that this
drift predicts the failure of naive transfer.

\subsection{Representational Alignment}

We restrict the transformation class $\mathcal{T}$ to act \emph{per module}
through a pair of linear maps that re-express the adapter in $M_{v+1}$'s
activation geometry. Let
$P^{\mathrm{in}}_\ell \in \mathbb{R}^{d^{\mathrm{in}}_\ell \times d^{\mathrm{in}}_\ell}$
map an input activation in $M_{v+1}$'s space to the corresponding point in
$M_v$'s space, and
$P^{\mathrm{out}}_\ell \in \mathbb{R}^{d^{\mathrm{out}}_\ell \times d^{\mathrm{out}}_\ell}$
map an output activation in $M_v$'s space to $M_{v+1}$'s space. The transplanted
update is
\begin{equation}
  \Delta W'_\ell
  \;=\; P^{\mathrm{out}}_\ell \, \Delta W_\ell \, P^{\mathrm{in}}_\ell
  \;=\; \tfrac{\alpha}{r}
        \underbrace{\left(P^{\mathrm{out}}_\ell B_\ell\right)}_{B'_\ell}
        \underbrace{\left(A_\ell P^{\mathrm{in}}_\ell\right)}_{A'_\ell},
\label{eq:transplant}
\end{equation}
so $T(A)_\ell = (B'_\ell, A'_\ell)$ \emph{preserves the rank-$r$ LoRA
structure}: portability costs no extra inference parameters. Intuitively,
$A'_\ell$ first pulls an $M_{v+1}$ input back into $M_v$'s basis, the original
adapter operates there, and $P^{\mathrm{out}}_\ell$ pushes the result forward
into $M_{v+1}$'s output basis.

\paragraph{Fitting the maps.}
The maps are estimated from anchor activations only. Passing
$\mathcal{D}_{\mathrm{anchor}}$ through both checkpoints yields paired
input-activation matrices $H^{v}_\ell, H^{v+1}_\ell$ (and analogously output
matrices). We solve a ridge-regularized least-squares problem
\begin{equation}
  P^{\mathrm{in}}_\ell
  \;=\;
  \arg\min_{P}\;
  \big\lVert H^{v}_\ell - H^{v+1}_\ell P^\top \big\rVert_F^2
  + \lambda \lVert P \rVert_F^2 ,
\label{eq:procrustes}
\end{equation}
optionally constraining $P$ to be orthogonal (the orthogonal Procrustes
solution $P = UV^\top$ from the SVD $H_v^\top H_{v+1} = U\Sigma V^\top$), which
encodes the prior that version updates are approximately rotations of a shared
feature space. $P^{\mathrm{out}}_\ell$ is fit symmetrically with the roles of
$v$ and $v{+}1$ swapped. Both are closed-form, require a single forward pass of
$\mathcal{D}_{\mathrm{anchor}}$ through each base model, and never touch
$\mathcal{D}_{\mathrm{task}}$.

\subsection{Assumptions and Scope}
\label{sec:assumptions}
We assume (i) $M_v$ and $M_{v+1}$ share architecture and adapted-module
dimensions (the common case for minor version bumps within a family), and
(ii) a shared or alignable tokenizer so that $\mathcal{D}_{\mathrm{anchor}}$
induces token-aligned activations. We treat differing tokenizers, differing
hidden dimensions, and cross-family transfer as out of scope for this paper and
return to them in \S\ref{sec:limitations}.
